\section{Implementierung der Software}
Zunächst werden grundlegende Funktionen zur Initialisierung der Hardware sowie zur Ansteuerung der LED-Streifen implementiert.  
Für die Ansteuerung der LEDs wird die Bibliothek \texttt{Adafruit\_NeoPixel} verwendet.  
Dabei werden für die Hauptanzeige, die Sekundenanzeige und die Anzeige der Zusatzfunktionen separate Objekte angelegt (vgl. Skript~\ref{lst:ledInit}).

\begin{lstlisting}[style=arduino,caption={Initialisierung der LED-Stränge.},label={lst:ledInit}]
#define LED_PIN            1
#define LED_COUNT          256

#define LED_PIN_SEC        7
#define LED_COUNT_SEC      60

#define LED_PIN_THIRD      10
#define LED_COUNT_THIRD    38

Adafruit_NeoPixel strip(LED_COUNT, LED_PIN, NEO_GRB + NEO_KHZ800);
Adafruit_NeoPixel secondStrip(LED_COUNT_SEC, LED_PIN_SEC, NEO_GRB + NEO_KHZ800);
Adafruit_NeoPixel thirdStrip(LED_COUNT_THIRD, LED_PIN_THIRD, NEO_GRB + NEO_KHZ800);
\end{lstlisting}

Eine Hilfsfunktion sorgt für die Aktivierung einzelner LEDs.  
Vor dem Setzen eines Farbwertes wird geprüft, ob sich der LED-Index innerhalb des gültigen Bereichs befindet (vgl. Skript~\ref{lst:IndexCheck}).

\begin{lstlisting}[style=arduino,caption={Setzen einer einzelnen LED.},label={lst:IndexCheck}]
void setLed(int ledNumber, uint8_t r, uint8_t g, uint8_t b) {
  if (ledNumber >= 0 && ledNumber < LED_COUNT) {
    strip.setPixelColor(ledNumber, strip.Color(r, g, b));
  }
}
\end{lstlisting}

Die Indexnummern der einzelnen \gls{led}s werden zur einfacheren Verwendung im Code globalen Variablen zugeordnet, welche die damit anzuzeigenden Wörter repräsentieren. 
Als Beispiel für die Uhrzeitanzeige dient die Variable „\texttt{HOUR\_2}“, welche das Wort „ZWEI“ auf der WordClock aktiviert und durch die \gls{led}s 71 bis 74 dargestellt wird (vgl. Skript~\ref{lst:WortSetzen}).
Die Verarbeitung der Zeit erfolgt durch zyklisches Auslesen der lokalen Systemzeit.  
Die zentrale Anzeigefunktion löscht zunächst den Hauptstreifen und aktiviert anschließend die benötigten Wörter für Stunden-, Minuten- und Zusatzinformationen (vgl. Skript~\ref{lst:WortSetzen}).

\begin{lstlisting}[style=arduino,caption={Hauptanzeigefunktion der WordClock.},label={lst:WortSetzen}]
int HOUR_2[]  = {71, 72, 73, 74};

void showHourWord(int hour24) {
  strip.clear();

  if (!lampPower) {
    strip.show();
    return;
  }

  showAlwaysOnLeds();

  int hourToShow = getDisplayHour(hour24, currentMinute);

  showMinuteWords(currentMinute);
  showHourNumber(hourToShow, currentMinute);
  showWeatherStatusLeds();

  strip.show();
}
\end{lstlisting}  

Zur Bedienung und Überwachung des Systems wird ein HTTP-Server auf dem ESP32 gestartet.
Die Weboberfläche zeigt aktuelle Statuswerte an und erlaubt die Einstellung von Ein-/Aus-Zustand, Helligkeit und Farbe (vgl. Skript~\ref{lst:HTTP}).
Die Helligkeitssteuerung erfolgt über eine Skalierung der \gls{rgb}-Werte mit der Funktion \texttt{applyBrightness} (vgl. Anhang~\ref{apx:1}).

\begin{lstlisting}[style=arduino,caption={Registrierung der Webserver-Routen.},label={lst:HTTP}]
void setupHttpServer() {
  server.on("/", handleRoot);
  server.on("/api/status", handleApiStatus);
  server.on("/set", handleSet);
  server.on("/weather", handleWeatherNow);
  server.onNotFound(handleNotFound);

  server.begin();

  Serial.println("HTTP-Server gestartet.");
  Serial.print("Im Browser oeffnen: http://");
  Serial.println(WiFi.localIP());
}
\end{lstlisting}

Die Hauptfunktionen sind \texttt{setup()} und \texttt{loop()}.
In der Funktion \texttt{setup()} werden Hardware, Sensorik, WLAN, Zeitsynchronisation, Wetterdaten und Webserver initialisiert (vgl. Skript~\ref{lst:Setup}).

\begin{lstlisting}[style=arduino,caption={Initialisierung der Programmlogik.},label={lst:Setup}]
void setup() {
  Serial.begin(115200);

  onboardLed.begin();
  onboardLed.setPixelColor(0, 0, 0, 0);
  onboardLed.show();

  strip.begin();
  strip.setBrightness(STRIP_BRIGHTNESS);
  strip.clear();
  strip.show();

  secondStrip.begin();
  secondStrip.setBrightness(STRIP_BRIGHTNESS);
  secondStrip.clear();
  secondStrip.show();

  thirdStrip.begin();
  thirdStrip.setBrightness(STRIP_BRIGHTNESS);
  thirdStrip.clear();
  thirdStrip.show();
  dht.begin();

  connectToWiFi();
  configTzTime(TZ_INFO, NTP_SERVER1, NTP_SERVER2);
  syncTimeUntilSuccessAtStartup();
  updateWeatherFromServer();
  updateIndoorSensor();

  bootModeOff();

  if (updateTimeVariables()) {
    showHourWord(currentHour);
  }

  setupHttpServer();
  Serial.println("System gestartet.");
}
\end{lstlisting}

In der Funktion \texttt{Loop()} werden unter anderem die Webserver-Anfragen verarbeitet. Die Funktion wird zyklisch wiederholt um alle Anzeigen zu aktualisieren wenn es eine Veränderung der Werte gibt (vgl. Skript~\ref{lst:Loop} und ~\ref{lst:Setup}).

\begin{lstlisting}[style=arduino,caption={Zyklischer Programmablauf.},label={lst:Loop}]
void loop() {
  server.handleClient();

  static unsigned long lastClockUpdate = 0;
  static unsigned long lastTimePrint = 0;

  unsigned long nowMs = millis();

  if (nowMs - lastClockUpdate >= 500) {
    lastClockUpdate = nowMs;

    if (updateTimeVariables()) {
      checkDailyResync();
      checkPeriodicWeatherUpdate();
      checkPeriodicIndoorSensorUpdate();

      showHourWord(currentHour);
      showSecondLed(currentSecond);
      showThirdStripStatus();
    }
  }

  if (nowMs - lastTimePrint >= 2000) {
    lastTimePrint = nowMs;

    Serial.printf("Aktuelle Zeit: %02d:%02d:%02d | RGB=(%d,%d,%d) | On=%d | Brightness=%d | Temp=%.1f | Code=%d | WetterValid=%d\n",
                  currentHour, currentMinute, currentSecond,
                  ledR, ledG, ledB, lampPower, webBrightness,
                  outsideTemperature, weatherCode, weatherValid);
  }
}
\end{lstlisting}
Die Software wird schrittweise erweitert, aufgespielt und getestet.  
Dieses Vorgehen ermöglicht eine schrittweise Integration aller Systemfunktionen und vereinfacht die Fehlersuche.
Der vollständige Programmcode befindet sich in Anhang~\ref{apx:1}.